\documentclass[12pt, a4paper]{article}
% ============================================================
%  preamble.tex — تنظیمات مشترک اسناد LaTeX (کامپایل با XeLaTeX)
% ============================================================
\usepackage{xcolor}
\usepackage{geometry}
\geometry{margin=2.5cm}
\usepackage{titlesec}
\usepackage{enumitem}
\setlist{nosep, leftmargin=*}
\usepackage{longtable}
\usepackage{booktabs}
\usepackage{array}
\usepackage{calc}
\usepackage{ragged2e}
\usepackage{fancyhdr}
\usepackage{etoolbox}
\usepackage{footnote}

% ---------- لازمهٔ جدول‌های خروجی pandoc ----------
\makeatletter
\patchcmd\longtable{\par}{\if@noskipsec\mbox{}\fi\par}{}{}
\makeatother
\makesavenoteenv{longtable}
\providecommand{\tightlist}{\setlength{\itemsep}{0pt}\setlength{\parskip}{0pt}}

% ---------- ظاهر ----------
\definecolor{primary}{RGB}{31, 78, 121}
\definecolor{accent}{RGB}{0, 150, 136}
\titleformat{\section}{\Large\bfseries\color{primary}}{\thesection}{1em}{}
\titleformat{\subsection}{\large\bfseries}{\thesubsection}{1em}{}
\titleformat{\subsubsection}{\normalsize\bfseries}{\thesubsubsection}{1em}{}
\renewcommand{\arraystretch}{1.3}
\linespread{1.3}

\pagestyle{fancy}
\fancyhf{}
\fancyfoot[C]{\thepage}
\renewcommand{\headrulewidth}{0.4pt}
\setlength{\headheight}{14pt}

% ---------- hyperref قبل از xepersian ----------
\usepackage{hyperref}
\hypersetup{colorlinks=true, linkcolor=primary, urlcolor=primary, citecolor=primary}

% ---------- فارسی‌سازی (آخرین بسته) ----------
\usepackage{xepersian}
\settextfont{Tahoma}
\setlatintextfont{Times New Roman}
\setdigitfont{Tahoma}

% ---------- شمارندهٔ لازم برای جدول‌های pandoc (LTcaptype = none) ----------
\newcounter{none}

\fancyhead[R]{\small سند پرسوناها — کلاس‌یار}
\fancyhead[L]{\small نسخهٔ ۱.۰}

\begin{document}

\begin{titlepage}
\centering
\vspace*{3cm}
{\Huge\bfseries سند پرسوناها\par}
\vspace{0.5cm}
{\LARGE\bfseries پروژهٔ کلاس‌یار\par}
\vspace{0.5cm}
{\large پلتفرم اشتراک جزوه، نمونه‌سؤال و ارزیابی استادان\par}
\vspace{2cm}
{\large تهیه‌کنندگان:\par}
{\large مریم قندی زنجانی، ملیکا شهبازی، آنیسا فروغمند، مائده رضایی، محمدمهدی ذوالقدر\par}
\vspace{0.5cm}
{\large دانشگاه آزاد زنجان\par}
\vspace{1cm}
{\large ۱۴۰۵/۰۷/۱۱\par}
\end{titlepage}

\tableofcontents
\newpage

{\def\LTcaptype{none} % do not increment counter
\begin{longtable}[]{@{}
  >{\raggedright\arraybackslash}p{(\linewidth - 2\tabcolsep) * \real{0.5000}}
  >{\raggedright\arraybackslash}p{(\linewidth - 2\tabcolsep) * \real{0.5000}}@{}}
\toprule\noalign{}
\endhead
\bottomrule\noalign{}
\endlastfoot
\textbf{نام پروژه} & کلاس‌یار (ClassYar) \\
\textbf{نسخهٔ سند} & ۱.۰ \\
\textbf{تاریخ} & ۱۴۰۵ \\
\textbf{تهیه‌کننده} & {[}مریم قندی زنجانی، ملیکا شهبازی،آنیسا فروغمند، مائده رضایی، محمدمهدی ذوالقدر {]} \\
\textbf{وضعیت} & پیش‌نویس \\
\end{longtable}
}

\begin{center}\rule{0.5\linewidth}{0.5pt}\end{center}

\section{مقدمه}

\subsection{هدف سند}

این سند پرسوناهای کاربران هدف پروژهٔ \textbf{کلاس‌یار} را تعریف می‌کند.
هدف آن ایجاد درک عمیق از نیازها، اهداف و دردهای کاربران است تا
تصمیم‌گیری‌های طراحی و توسعه بر پایهٔ آن‌ها انجام شود.

\subsection{چرا پرسونا؟}

پرسونا یک \textbf{نمایندهٔ فرضی} از یک گروه کاربری است که با داده‌های واقعی
ساخته می‌شود. پرسونا کمک می‌کند:
- تیم توسعه بداند برای \textbf{چه کسی} می‌سازد.
- از «طراحی برای همه» پرهیز شود.
- اولویت‌بندی ویژگی‌ها بر پایهٔ نیاز واقعی کاربران باشد.

\subsection{روش‌شناسی}

پرسوناها بر پایهٔ موارد زیر ساخته شده‌اند:
- مصاحبهٔ فرضی با دانشجویان و یک استاد
- مشاهدهٔ گروه‌های کلاسی در پیام‌رسان‌ها
- بررسی سایت‌های مشابه اشتراک جزوه
- تجربهٔ شخصی تیم به‌عنوان دانشجو

\begin{center}\rule{0.5\linewidth}{0.5pt}\end{center}

\section{پرسونا ۱ — نگار (دانشجوی استفاده‌کننده)}

\subsection{مشخصات دموگرافیک}

{\def\LTcaptype{none} % do not increment counter
\begin{longtable}[]{@{}ll@{}}
\toprule\noalign{}
ویژگی & مقدار \\
\midrule\noalign{}
\endhead
\bottomrule\noalign{}
\endlastfoot
نام & نگار حیدری \\
سن & ۲۰ سال \\
جنسیت & زن \\
شغل & دانشجوی کارشناسی مهندسی کامپیوتر \\
محل زندگی & رشت \\
تحصیلات & کارشناسی (ترم ۳) \\
مهارت فنی & بالا \\
وضعیت تأهل & مجرد \\
\end{longtable}
}

\subsection{نقل‌قول}

\begin{quote}
«شب امتحان بین ده تا گروه دنبال جزوهٔ این استاد می‌گردم و نمی‌دانم کدام‌یک کامل است.»
\end{quote}

\subsection{اهداف}

\begin{itemize}
\tightlist
\item
  پیدا کردن سریع جزوهٔ استاد و درس موردنظر
\item
  اطمینان از کیفیت جزوه قبل از دانلود
\item
  دسترسی به نمونه‌سؤال‌های همان استاد
\item
  شناخت استاد قبل از انتخاب واحد
\end{itemize}

\subsection{دردها}

\begin{itemize}
\tightlist
\item
  جزوه‌ها در گروه‌های مختلف پخش شده‌اند
\item
  معلوم نیست جزوه کامل و درست است یا نه
\item
  نمی‌تواند بر پایهٔ استاد و دانشگاه جست‌وجو کند
\item
  دربارهٔ استاد اطلاعات مطمئنی ندارد
\end{itemize}

\subsection{نیازها}

\begin{itemize}
\tightlist
\item
  فیلتر بر پایهٔ استاد، استان، واحد و درس
\item
  دیدن امتیاز و نظر دیگران
\item
  دانلود ساده
\item
  دیدن نمونه‌سؤال و لینک گروه کلاس
\end{itemize}

\subsection{سناریوی استفاده}

\begin{quote}
شب قبل از امتحان با موبایل وارد سایت می‌شود. استاد و درس را انتخاب می‌کند
و واحد دانشگاه خودش را فیلتر می‌کند. سه جزوه می‌بیند، امتیاز و نظرها را می‌خواند،
بهترین را دانلود می‌کند و بعد از امتحان برایش امتیاز ثبت می‌کند.
\end{quote}

\begin{center}\rule{0.5\linewidth}{0.5pt}\end{center}

\section{پرسونا ۲ — امیرحسین (دانشجوی مشارکت‌کننده)}

\subsection{مشخصات دموگرافیک}

{\def\LTcaptype{none} % do not increment counter
\begin{longtable}[]{@{}ll@{}}
\toprule\noalign{}
ویژگی & مقدار \\
\midrule\noalign{}
\endhead
\bottomrule\noalign{}
\endlastfoot
نام & امیرحسین کریمی \\
سن & ۲۲ سال \\
جنسیت & مرد \\
شغل & دانشجوی کارشناسی مهندسی صنایع \\
محل زندگی & تبریز \\
تحصیلات & کارشناسی (ترم ۶) \\
مهارت فنی & متوسط \\
وضعیت تأهل & مجرد \\
\end{longtable}
}

\subsection{نقل‌قول}

\begin{quote}
«جزوه‌هایم کامل است و دوست دارم به درد بقیه هم بخورد، به شرطی که دیده شود.»
\end{quote}

\subsection{اهداف}

\begin{itemize}
\tightlist
\item
  اشتراک جزوه‌های دست‌نویس خودش
\item
  گرفتن بازخورد و امتیاز
\item
  دیده‌شدن کارش توسط استاد
\item
  گرفتن نمرهٔ تشویقی
\end{itemize}

\subsection{دردها}

\begin{itemize}
\tightlist
\item
  ارسال فایل در گروه‌ها زود گم می‌شود
\item
  هیچ بازخوردی دربارهٔ کیفیت جزوه نمی‌گیرد
\item
  زحمتش دیده نمی‌شود
\item
  اسکن و ارسال تکراری وقت‌گیر است
\end{itemize}

\subsection{نیازها}

\begin{itemize}
\tightlist
\item
  آپلود ساده با انتخاب استاد، درس و ترم
\item
  دیدن امتیاز و نظر دربارهٔ جزوه
\item
  مدیریت و ویرایش جزوه‌های خودش
\item
  اطلاع از تأیید استاد
\end{itemize}

\subsection{سناریوی استفاده}

\begin{quote}
جزوهٔ اسکن‌شدهٔ «تحقیق در عملیات» را انتخاب می‌کند. استاد، درس و ترم را
وارد می‌کند و آپلود می‌کند. چند روز بعد می‌بیند ۱۲ نفر امتیاز داده‌اند
و استاد آن را تأیید کرده است.
\end{quote}

\begin{center}\rule{0.5\linewidth}{0.5pt}\end{center}

\section{پرسونا ۳ — دکتر کاظمی (استاد)}

\subsection{مشخصات دموگرافیک}

{\def\LTcaptype{none} % do not increment counter
\begin{longtable}[]{@{}ll@{}}
\toprule\noalign{}
ویژگی & مقدار \\
\midrule\noalign{}
\endhead
\bottomrule\noalign{}
\endlastfoot
نام & دکتر مهدی کاظمی \\
سن & ۴۶ سال \\
جنسیت & مرد \\
شغل & عضو هیئت علمی دانشگاه \\
محل زندگی & اصفهان \\
تحصیلات & دکتری \\
مهارت فنی & متوسط \\
وضعیت تأهل & متأهل، دو فرزند \\
\end{longtable}
}

\subsection{نقل‌قول}

\begin{quote}
«دوست دارم بدانم دانشجوها از چه منابعی درس می‌خوانند و فعال‌ها را تشویق کنم.»
\end{quote}

\subsection{اهداف}

\begin{itemize}
\tightlist
\item
  دیدن منابع دانشجویی مربوط به درس خودش
\item
  تأیید جزوه‌های درست و کامل
\item
  تشویق دانشجویان فعال با نمرهٔ تشویقی
\item
  شنیدن بازخورد دانشجویان
\end{itemize}

\subsection{دردها}

\begin{itemize}
\tightlist
\item
  از منابع غیررسمی دانشجویان خبر ندارد
\item
  ممکن است جزوهٔ اشتباه بین دانشجویان پخش شود
\item
  راهی برای تشویق دانشجویان فعال ندارد
\item
  بازخورد صادقانه کم دریافت می‌کند
\end{itemize}

\subsection{نیازها}

\begin{itemize}
\tightlist
\item
  پنل ساده و فارسی
\item
  فهرست جزوه‌های مربوط به درس خودش
\item
  تأیید جزوه و ثبت نمرهٔ تشویقی با چند کلیک
\item
  امکان پاسخ به نظرها
\end{itemize}

\subsection{سناریوی استفاده}

\begin{quote}
با لپ‌تاپ وارد سامانه می‌شود. فهرست جزوه‌های درس «تحقیق در عملیات» را می‌بیند.
یکی را بررسی و تأیید می‌کند و برای آپلودکننده نمرهٔ تشویقی ثبت می‌کند.
زیر یک نظر دربارهٔ خودش پاسخ می‌دهد.
\end{quote}

\begin{center}\rule{0.5\linewidth}{0.5pt}\end{center}

\section{پرسونا ۴ — خانم نوروزی (مدیر سامانه)}

\subsection{مشخصات دموگرافیک}

{\def\LTcaptype{none} % do not increment counter
\begin{longtable}[]{@{}ll@{}}
\toprule\noalign{}
ویژگی & مقدار \\
\midrule\noalign{}
\endhead
\bottomrule\noalign{}
\endlastfoot
نام & سمیرا نوروزی \\
سن & ۲۶ سال \\
جنسیت & زن \\
شغل & دانشجوی ارشد و مسئول انجمن علمی دانشکده \\
محل زندگی & تهران \\
تحصیلات & کارشناسی ارشد \\
مهارت فنی & بالا \\
وضعیت تأهل & مجرد \\
\end{longtable}
}

\subsection{نقل‌قول}

\begin{quote}
«اگر کنترلی روی محتوا نباشد، اعتماد کاربران از بین می‌رود.»
\end{quote}

\subsection{اهداف}

\begin{itemize}
\tightlist
\item
  حفظ نظم و اعتبار محتوا
\item
  مدیریت استادان، دانشگاه‌ها و واحدها
\item
  رسیدگی سریع به گزارش‌ها
\item
  رشد تعداد کاربران
\end{itemize}

\subsection{دردها}

\begin{itemize}
\tightlist
\item
  ممکن است محتوای نامناسب یا دارای حق مالکیت منتشر شود
\item
  نظرهای توهین‌آمیز دربارهٔ استادان
\item
  مدیریت دستی اطلاعات پایه وقت‌گیر است
\item
  گزارش‌ها پراکنده‌اند
\end{itemize}

\subsection{نیازها}

\begin{itemize}
\tightlist
\item
  پنل مدیریت ساده
\item
  فهرست گزارش‌ها در یک صفحه
\item
  حذف یا پنهان‌کردن محتوا و نظر
\item
  افزودن استان، دانشگاه، واحد و استاد
\end{itemize}

\subsection{سناریوی استفاده}

\begin{quote}
صبح وارد پنل مدیریت می‌شود. دو گزارش جدید می‌بیند. یک نظر توهین‌آمیز را حذف می‌کند،
گزارش دیگر را رد می‌کند و درخواست افزودن یک استاد جدید را تأیید می‌کند.
\end{quote}

\begin{center}\rule{0.5\linewidth}{0.5pt}\end{center}

\section{ماتریس اولویت‌بندی}

{\def\LTcaptype{none} % do not increment counter
\begin{longtable}[]{@{}lll@{}}
\toprule\noalign{}
پرسونا & تأثیر & اولویت \\
\midrule\noalign{}
\endhead
\bottomrule\noalign{}
\endlastfoot
نگار (دانشجوی استفاده‌کننده) & بالا & اول \\
امیرحسین (دانشجوی مشارکت‌کننده) & بالا & اول \\
دکتر کاظمی (استاد) & بالا & دوم \\
خانم نوروزی (مدیر سامانه) & متوسط & دوم \\
\end{longtable}
}

\begin{center}\rule{0.5\linewidth}{0.5pt}\end{center}

\section*{پیوست الف — تاریخچهٔ تغییرات}\addcontentsline{toc}{section}{پیوست الف — تاریخچهٔ تغییرات}

{\def\LTcaptype{none} % do not increment counter
\begin{longtable}[]{@{}
  >{\raggedright\arraybackslash}p{(\linewidth - 6\tabcolsep) * \real{0.1842}}
  >{\raggedright\arraybackslash}p{(\linewidth - 6\tabcolsep) * \real{0.1842}}
  >{\raggedright\arraybackslash}p{(\linewidth - 6\tabcolsep) * \real{0.2632}}
  >{\raggedright\arraybackslash}p{(\linewidth - 6\tabcolsep) * \real{0.3684}}@{}}
\toprule\noalign{}
\begin{minipage}[b]{\linewidth}\raggedright
نسخه
\end{minipage} & \begin{minipage}[b]{\linewidth}\raggedright
تاریخ
\end{minipage} & \begin{minipage}[b]{\linewidth}\raggedright
تغییرات
\end{minipage} & \begin{minipage}[b]{\linewidth}\raggedright
تهیه‌کننده
\end{minipage} \\
\midrule\noalign{}
\endhead
\bottomrule\noalign{}
\endlastfoot
۱.۰ & ۱۴۰۵ & نسخهٔ اولیه & {[}مریم قندی زنجانی، ملیکا شهبازی،آنیسا فروغمند، مائده رضایی، محمدمهدی ذوالقدر{]} \\
\end{longtable}
}

\end{document}
